\chapter{Implementación}

\section{Estructura del repositorio}

La implementación se organiza en tres áreas:

\begin{itemize}
  \item `codigo/`: aplicación, tests, datos, modelos, informes y documentación
        técnica;
  \item `memoria/`: documento LaTeX del TFM;
  \item `recursos/`: PDFs, referencias y diagramas auxiliares.
\end{itemize}

\section{Estado Actual}

Hasta este punto se ha creado la estructura inicial de directorios, se han
descargado los datasets crudos, se ha definido el flujo de datos del MVP y se
han implementado el contrato inicial del estado global del grafo y los esquemas
Pydantic de entradas, decisiones y resultados. Los datos crudos se conservan en
`codigo/data/raw/`, junto con un README de trazabilidad.

\section{Estado global implementado}

El estado global se ha dividido en dos capas:

\begin{itemize}
  \item `codigo/app/schemas/state.py`: modelos Pydantic estrictos para validar
        el estado, configuraciones, métricas, errores y artefactos;
  \item `codigo/app/graph/state.py`: `TypedDict` usado por LangGraph,
        validador de estado e inicializador del caso CWRU.
\end{itemize}

El inicializador `create\_initial\_cwru\_state` crea una ejecución preparada para
generar el manifiesto del dataset. El estado resultante no contiene señales ni
arrays, solo rutas, contexto experimental y campos reservados para artefactos
futuros.

\section{Contratos Pydantic implementados}

Se han añadido tres grupos de contratos:

\begin{itemize}
  \item `codigo/app/schemas/dataset.py`: manifiesto del dataset, canales y
        metadatos de fallo;
  \item `codigo/app/schemas/agent\_decisions.py`: decisiones estructuradas del
        supervisor, limpiador, estructurador, modelador, evaluador y redactor;
  \item `codigo/app/schemas/executor\_results.py`: resultados comunes y
        especializados de ejecutores deterministas.
\end{itemize}

Estos contratos consolidan la regla arquitectónica de que los agentes producen
configuraciones o decisiones validables, mientras que los ejecutores devuelven
rutas, métricas, errores y referencias a artefactos, pero no datos pesados.

\section{Artefactos previstos}

\begin{itemize}
  \item `codigo/data/interim/cwru\_bearing/manifest.csv`;
  \item `codigo/data/interim/cwru\_bearing/profile.json`;
  \item `codigo/data/processed/cwru\_bearing/clean\_signals/`;
  \item `codigo/data/tensors/cwru\_bearing/windows\_features.csv`;
  \item `codigo/data/tensors/cwru\_bearing/windows\_raw.npz`;
  \item `codigo/data/tensors/cwru\_bearing/splits.json`;
  \item `codigo/reports/cwru\_bearing/`;
  \item `codigo/models/cwru\_bearing/`.
\end{itemize}

\section{Adaptadores de Entrada}

Para reducir el acoplamiento al formato CWRU se ha introducido una capa ligera
de adaptadores en `codigo/app/services/signal\_adapters.py`. Esta capa expone
dos operaciones comunes: leer todos los canales numéricos de un fichero y cargar
un canal concreto como vector. Los ejecutores de perfilado y limpieza usan ahora
esta frontera en lugar de leer directamente ficheros `.mat`.

La implementación inicial soporta `.mat`, `.csv`, `.txt`, `.tsv` y `.npz`. En
el caso de CWRU, el adaptador reconoce los canales `DE\_time`, `FE\_time`,
`BA\_time` y `RPM`; en formatos tabulares, interpreta columnas numéricas como
canales. Para incorporar un nuevo origen, como NASA IMS, el objetivo será crear
o ampliar un adaptador de entrada y mantener estable el resto del pipeline:
perfilado, limpieza, estructuración, modelado y evaluación.

\section{Ejecutor de manifiesto}

Se ha implementado el primer ejecutor determinista del pipeline, el módulo
`dataset\_manifest.py`. Su responsabilidad es generar el artefacto:

\begin{center}
\begin{minipage}{0.85\textwidth}
\begin{verbatim}
codigo/data/interim/cwru_bearing/manifest.csv
\end{verbatim}
\end{minipage}
\end{center}

a partir de los ficheros `.mat` crudos de CWRU.

El ejecutor no carga señales completas. Solo inspecciona los nombres de fichero
y aplica un mapeo controlado de metadatos CWRU: carga, RPM, tipo de fallo,
diámetro y posición del fallo exterior cuando corresponde. El manifiesto
generado conserva la frecuencia original de cada fichero y fija una frecuencia
objetivo de 12 kHz para el MVP.

La ejecución real sobre los datos descargados ha generado 64 registros: 4
ficheros normales y 60 ficheros con fallo. El fichero resultante contiene 65
líneas incluyendo la cabecera.

\section{Ejecutor de perfilado}

Se ha implementado el perfilador determinista `data\_profiler.py`. Este ejecutor
lee el manifiesto, inspecciona los ficheros `.mat` y genera:

\begin{center}
\begin{minipage}{0.85\textwidth}
\begin{verbatim}
codigo/data/interim/cwru_bearing/profile.json
\end{verbatim}
\end{minipage}
\end{center}

El perfil contiene metadatos y estadísticas por canal: número de muestras,
valores no finitos, mínimo, máximo, media, desviación típica, RMS, curtosis y
energía. No almacena arrays ni señales completas.

La ejecución real ha perfilado 64 ficheros. El resumen agregado contiene 4
ficheros normales, 60 ficheros con fallo, frecuencias de muestreo de 48 kHz y
12 kHz, y los canales detectados `BA\_time`, `DE\_time`, `FE\_time` y `RPM`.

\section{Ejecutor de Limpieza}

Se ha implementado el ejecutor determinista `cleaning.py`. Este ejecutor lee el
manifiesto y el perfil, extrae el canal principal configurado, elimina valores
no finitos, remuestrea la señal a 12 kHz cuando procede y guarda una señal
limpia comprimida por fichero:

\begin{center}
\begin{minipage}{0.85\textwidth}
\begin{verbatim}
codigo/data/processed/cwru_bearing/clean_signals/
\end{verbatim}
\end{minipage}
\end{center}

También genera un resumen de auditoría en:

\begin{center}
\begin{minipage}{0.85\textwidth}
\begin{verbatim}
codigo/data/processed/cwru_bearing/cleaning_summary.json
\end{verbatim}
\end{minipage}
\end{center}

La ejecución real ha producido 64 ficheros `.npz` comprimidos. Los cuatro
ficheros normales se han remuestreado de 48 kHz a 12 kHz; los ficheros de fallo
ya estaban a 12 kHz. El estado del grafo conserva únicamente la ruta del
directorio limpio y las referencias a artefactos, no las señales.

\section{Ejecutor de Estructuración Temporal}

Se ha implementado el ejecutor determinista `structuring.py`. Este ejecutor lee
las señales limpias, valida que el canal y la frecuencia coincidan con la
configuración, genera ventanas temporales con tamaño 2048 y solapamiento del
50\%, calcula features temporales iniciales y crea particiones reproducibles a
nivel de fichero.

\begin{center}
\begin{minipage}{0.85\textwidth}
\begin{verbatim}
codigo/data/tensors/cwru_bearing/windows_features.csv
codigo/data/tensors/cwru_bearing/windows_raw.npz
codigo/data/tensors/cwru_bearing/splits.json
\end{verbatim}
\end{minipage}
\end{center}

La estrategia de partición evita mezclar ventanas de un mismo fichero entre
subconjuntos. Las ventanas normales se reparten entre entrenamiento,
validación y test, mientras que las ventanas con fallo se reservan para test.
La ejecución real ha generado 7460 ventanas de 2048 muestras: 175 de
entrenamiento, 117 de validación y 7168 de test. El tensor crudo tiene forma
`(7460, 2048)` y el directorio de tensores ocupa aproximadamente 22 MB.

\section{Ejecutor de Modelado Base}

Se ha implementado el ejecutor determinista `modeling.py` con el primer modelo
base de detección de anomalías: Isolation Forest. El modelo se entrena
únicamente con ventanas normales del subconjunto de entrenamiento y consume las
features temporales generadas por la fase anterior.

\begin{center}
\begin{minipage}{0.85\textwidth}
\begin{verbatim}
codigo/models/cwru_bearing/isolation_forest.joblib
codigo/models/cwru_bearing/predictions.csv
codigo/models/cwru_bearing/modeling_summary.json
\end{verbatim}
\end{minipage}
\end{center}

El artefacto `joblib` conserva el modelo entrenado, las columnas usadas y el
umbral de decisión. Las predicciones guardan, por cada ventana, identificador,
fichero original, partición, etiqueta real, puntuación de anomalía, umbral y
predicción binaria. Las métricas finales no se calculan en este ejecutor, sino
que se reservan para la fase de evaluación.

La ejecución real ha entrenado el modelo con 175 ventanas normales y ha
generado 7460 predicciones. El umbral estimado ha sido 0.6334 y el resumen de
modelado registra 403 ventanas predichas como normales y 7057 como anómalas.
Estos conteos son una salida preliminar del modelo, no una evaluación final de
rendimiento.

\section{Ejecutor de Evaluación}

Se ha implementado el ejecutor determinista `evaluation.py`. Este módulo lee las
predicciones del modelo, compara la etiqueta binaria real con la predicción de
anomalía y calcula métricas técnicas por partición. La partición principal para
el MVP es `test`.

\begin{center}
\begin{minipage}{0.85\textwidth}
\begin{verbatim}
codigo/reports/cwru_bearing/evaluation/metrics.json
codigo/reports/cwru_bearing/evaluation/evaluation_summary.md
\end{verbatim}
\end{minipage}
\end{center}

La ejecución real sobre el split de test ha obtenido precision 0.9997, recall
1.0000, F1-score 0.9999, ROC-AUC 1.0000, PR-AUC 1.0000 y una tasa de falsos
positivos de 0.0171. La matriz de confusión ha sido: 115 verdaderos negativos,
2 falsos positivos, 0 falsos negativos y 7051 verdaderos positivos.

Estos resultados deben interpretarse como validación del pipeline y del
baseline inicial sobre un benchmark controlado. No constituyen todavía una
garantía de generalización industrial, que se deberá analizar con validaciones
más exigentes y, posteriormente, con otros datasets como NASA IMS.

\section{Dependencias}

Se ha creado `requirements.txt` como punto inicial para construir el entorno de
Miniconda `tfm\_v2`, fijado sobre Python 3.11. Incluye dependencias de contratos
y grafo, procesamiento de datos, modelos clásicos, visualización técnica y
verificación.

\section{Verificación Actual}

La verificación de los contratos se ha realizado con tests unitarios en el
entorno objetivo:

\begin{center}
\begin{minipage}{0.85\textwidth}
\begin{verbatim}
conda run -n tfm_v2 python -m unittest discover codigo/tests
\end{verbatim}
\end{minipage}
\end{center}

Las pruebas validan la creación del estado inicial CWRU, serialización JSON,
rechazo de campos no declarados, restricciones de configuraciones y métricas,
coherencia del manifiesto, decisiones de agentes, resultados de ejecutores y
generación del manifiesto, perfilado, limpieza, estructuración temporal y
modelado base y evaluación CWRU. También validan los adaptadores de entrada para
formatos `.mat` y `.csv`. En esta fase se han ejecutado 48 tests
correctamente.

\section{Siguiente Paso de Implementación}

El siguiente avance técnico será montar el grafo LangGraph mínimo, encadenando
los ejecutores deterministas y preparando los puntos donde después entrarán los
agentes especializados.
